\documentclass[11pt,twocolumn]{article}

% ── Packages ─────────────────────────────────────────────────────────────
\usepackage[utf8]{inputenc}
\usepackage[T1]{fontenc}
\usepackage{amsmath,amssymb}
\usepackage{booktabs}
\usepackage{hyperref}
\usepackage{graphicx}
\usepackage{xcolor}
\usepackage[margin=1in]{geometry}
\usepackage{caption}
\usepackage{enumitem}
\usepackage{listings}
\usepackage{authblk}

\lstset{
  language=Python,
  basicstyle=\ttfamily\scriptsize,
  keywordstyle=\bfseries\color{blue!70!black},
  commentstyle=\itshape\color{gray},
  breaklines=true,
  frame=single,
  framesep=2pt,
  columns=fullflexible,
}

\hypersetup{
  colorlinks=true,
  linkcolor=blue!70!black,
  citecolor=blue!70!black,
  urlcolor=blue!70!black,
}

% ── Title ────────────────────────────────────────────────────────────────
\title{Aethermor: An Open-Source Inverse Thermal Design Toolkit \\
       for Thermodynamic Computing Research}

\author{Aethermor Contributors}
\date{March 2026}

\begin{document}
\maketitle

% ── Abstract ─────────────────────────────────────────────────────────────
\begin{abstract}
As transistor scaling approaches fundamental thermodynamic limits,
hardware architects increasingly need tools that quantify how close a
design sits to the Landauer limit and what thermal constraints govern
achievable compute density. We present \textbf{Aethermor}, an
open-source Python toolkit that solves \emph{inverse} thermal design
problems: given a power budget, cooling architecture, and substrate
material, it finds the maximum gate density, minimum required cooling,
optimal power distribution across heterogeneous functional blocks, and
paradigm crossover points --- questions that currently require weeks of
manual sweeps in commercial finite-element tools. Aethermor combines a
3D Fourier thermal solver (0.00\% energy conservation error), ITRS/IRDS-calibrated CMOS
energy models, adiabatic and reversible computing paradigms, nine
substrate materials with CRC Handbook properties, multi-layer cooling
stack models, and heterogeneous chip floorplan simulation. A built-in
validation suite of 133 cross-checks against CODATA~2018, analytical
solutions, and published roadmap data ensures physical correctness. We
demonstrate that diamond substrates sustain $17\times$ higher gate
density than silicon at equal cooling, that adiabatic logic provides a
$2{,}006\times$ energy advantage over CMOS at 1~GHz on a 7~nm node,
and that thermal headroom redistribution on a heterogeneous SoC yields
up to $2\times$ throughput improvement. Aethermor is released under the
Apache~2.0 license with an interactive browser-based explorer, 254
automated tests, and seven ready-to-run research examples.
\end{abstract}

% ── 1. Introduction ─────────────────────────────────────────────────────
\section{Introduction}
\label{sec:intro}

The semiconductor industry faces a fundamental physical barrier:
Dennard scaling has ended, and each new technology node brings
diminishing energy-per-switch improvements while power density
continues to rise~\cite{dennard1974,esmaeilzadeh2011}. At advanced
nodes ($\leq 7$~nm), the energy per CMOS gate switch is approximately
$2.4 \times 10^{-16}$~J --- still roughly $85{,}000\times$ above the
Landauer limit of $k_B T \ln 2 \approx 2.87 \times 10^{-21}$~J at
room temperature~\cite{landauer1961}. This gap defines the
thermodynamic design space that hardware architects must navigate.

Existing thermal simulation tools fall into two categories.
\emph{Forward solvers} (HotSpot~\cite{skadron2004},
COMSOL~Multiphysics) compute temperature given a known power map.
\emph{Commercial FEM tools} can in principle solve inverse problems,
but require manual scripting, cost \$25{,}000+ per year, and lack
built-in models for Landauer-aware energy accounting, multi-paradigm
comparisons, and technology roadmap projection. No existing open-source
tool answers questions such as:
\begin{itemize}[nosep]
  \item What is the maximum gate density this substrate can sustain?
  \item What is the minimum cooling coefficient needed for this density?
  \item How should I distribute compute across functional blocks to
        maximise throughput under a thermal limit?
  \item At what frequency does adiabatic logic become more efficient
        than CMOS on this material?
\end{itemize}

Aethermor fills this gap. It is a Python toolkit that combines
SI-unit physics models with inverse design algorithms, enabling
hardware researchers to explore the full thermodynamic design space
interactively or programmatically. This paper describes the
architecture, physics models, validation methodology, and
representative results.


% ── 2. Architecture ─────────────────────────────────────────────────────
\section{System Architecture}
\label{sec:arch}

Aethermor is organized into three layers, each building on the one
below it.

\subsection{Physics Layer (\texttt{physics/})}

\paragraph{Fundamental Constants.}
All calculations derive from CODATA~2018 exact values: Boltzmann
$k_B = 1.380\,649 \times 10^{-23}$~J/K, Planck
$h = 6.626\,070\,15 \times 10^{-34}$~J$\cdot$s. The Landauer limit is
computed as $E_{\min}(T) = k_B \, T \, \ln 2$.

\paragraph{Material Database.}
Nine chip substrates are included (Table~\ref{tab:materials}), each
with published thermal conductivity, specific heat, density,
electrical resistivity, bandgap, and maximum operating temperature.
All values are sourced from the CRC Handbook of Chemistry and
Physics~\cite{crc97} and Morko\c{c}~\cite{morkoc2006,morkoc2008}
for wide-bandgap materials.

\begin{table}[t]
\centering
\caption{Substrate materials in Aethermor.}
\label{tab:materials}
\small
\begin{tabular}{@{}lrrr@{}}
\toprule
Material & $k$ (W/m$\cdot$K) & $T_{\max}$ (K) & $E_g$ (eV) \\
\midrule
Silicon          & 148   & 723  & 1.12 \\
Diamond          & 2{,}200 & 973  & 5.47 \\
SiC (4H)         & 490   & 873  & 3.26 \\
GaN              & 130   & 873  & 3.40 \\
GaAs             & 55    & 623  & 1.42 \\
SiO$_2$          & 1.4   & 1{,}673 & 9.0 \\
Copper           & 401   & 673  & ---  \\
InP              & 68    & 673  & 1.35 \\
Graphene         & 5{,}000 & 973  & 0.0  \\
\bottomrule
\end{tabular}
\end{table}

\paragraph{Gate Energy Models.}
Four computing paradigms are modelled:
\begin{enumerate}[nosep]
  \item \textbf{CMOS}: $E = \tfrac{1}{2} C_L V_{dd}^2 + E_{\mathrm{leak}}(T)$,
    with $V_{dd}$ and $C_L$ calibrated to ITRS~2013~\cite{itrs2013} and
    IRDS~2022~\cite{irds2022} roadmaps across 10 nodes (130--1.4~nm).
    Leakage follows an Arrhenius model: $I_{\mathrm{leak}}(T) = I_0
    \exp\!\bigl(\tfrac{T - T_0}{T_{\mathrm{coeff}}}\bigr)$.
  \item \textbf{Adiabatic}: $E = R C_L^2 V_{dd}^2 f$, proportional to
    frequency --- energy approaches zero as $f \to 0$.
  \item \textbf{Reversible}: $E = n_{\mathrm{erase}} \cdot k_B T \ln 2 +
    E_{\mathrm{overhead}}$, where $n_{\mathrm{erase}}$ is the number
    of irreversible bit erasures per gate operation.
  \item \textbf{Landauer floor}: $E = k_B T \ln 2$, the
    information-theoretic minimum.
\end{enumerate}

\paragraph{Fourier Thermal Solver.}
A 3D explicit finite-difference solver for the heat equation:
\begin{equation}
  \rho \, c_p \frac{\partial T}{\partial t}
    = k \, \nabla^2 T + \dot{q}
  \label{eq:fourier}
\end{equation}
with CFL-stable timestep $\Delta t < \tfrac{\Delta x^2}{6\alpha}$
($\alpha = k / \rho c_p$). The solver supports convective, fixed-temperature,
and adiabatic boundary conditions. Energy conservation is verified at every
run: $|E_{\mathrm{in}} - E_{\mathrm{out}} - \Delta U| / E_{\mathrm{in}} < 5\%$,
with observed error of 0.00\% in all validation runs.

\paragraph{Cooling Stack.}
A \texttt{CoolingStack} module models the die-to-ambient thermal path
as a series of thermal resistances:
\begin{equation}
  R_{\mathrm{total}} = \sum_i \frac{t_i}{k_i \cdot A} + \frac{1}{h_{\mathrm{amb}} \cdot A}
  \label{eq:rstack}
\end{equation}
Eleven pre-built layers (thermal paste, liquid metal, copper IHS,
aluminium heatsink, diamond spreader, etc.) and six factory
configurations (bare-die, desktop-air, server-air, liquid, direct-liquid,
diamond-spreader-liquid) are provided.

\paragraph{Chip Floorplan.}
A \texttt{ChipFloorplan} class defines heterogeneous architectures
with distinct functional blocks (CPU, GPU, cache, I/O), each with
independent gate density, activity factor, technology node, and
computing paradigm. Factory methods generate representative SoC and
hybrid CMOS/adiabatic layouts.

\subsection{Analysis Layer (\texttt{analysis/})}

The analysis layer provides eight inverse thermal design tools
(Section~\ref{sec:inverse}), technology roadmap projection, Landauer
gap analysis, design space sweeps with Pareto extraction, thermal
regime classification, and hotspot detection.

\subsection{Interface Layer}

Users interact with Aethermor through three channels: (1)~a
browser-based interactive explorer (\texttt{aethermor dashboard}) with six
parameterised tabs updated in real time; (2)~a Python scripting API;
(3)~seven ready-to-run example scripts demonstrating research workflows.


% ── 3. Inverse Thermal Design ───────────────────────────────────────────
\section{Inverse Thermal Design}
\label{sec:inverse}

The central contribution of Aethermor is solving the inverse thermal
design problem: \emph{given constraints, find the best design}. Eight
tools are provided, each addressing a distinct research question.

\subsection{Maximum Achievable Gate Density}

Given a material and cooling coefficient $h$, a binary search over the
3D Fourier solver finds the maximum gate density $D_{\max}$ such that
$T_{\max} \leq T_{\mathrm{limit}}$:
\begin{equation}
  D_{\max}: \quad T_{\max}\bigl(D_{\max}, k, h\bigr) = T_{\mathrm{limit}}
  \label{eq:dmax}
\end{equation}

\subsection{Minimum Required Cooling}

Given a material and target density $D$, the combined conduction +
convection 1D steady-state model:
\begin{equation}
  T_{\max} = T_{\mathrm{amb}} + Q \left( \frac{1}{h A} + \frac{\Delta x}{2 k A} \right)
  \label{eq:tmax1d}
\end{equation}
is inverted analytically for $h_{\min}$:
\begin{equation}
  h_{\min} = \frac{Q}{A \bigl( \Delta T - \Delta T_{\mathrm{cond}} \bigr)}
  \label{eq:hmin}
\end{equation}
where $\Delta T_{\mathrm{cond}} = Q \Delta x / (2kA)$ is the
irreducible conduction floor. If $\Delta T_{\mathrm{cond}} \geq
\Delta T$, no amount of convective cooling suffices --- the design
requires a higher-conductivity substrate.

\subsection{Material Ranking}

Runs \texttt{find\_max\_density} across all materials and sorts by
$D_{\max}$. At $h = 1{,}000$~W/(m$^2\cdot$K) on a 7~nm node,
diamond sustains $7.0 \times 10^8$ gates/element versus silicon's
$4.1 \times 10^7$ --- a $17\times$ advantage (Table~\ref{tab:ranking}).

\begin{table}[t]
\centering
\caption{Material ranking at 7~nm, 1~GHz, $h = 1{,}000$~W/(m$^2\cdot$K).}
\label{tab:ranking}
\small
\begin{tabular}{@{}lrl@{}}
\toprule
Material & $D_{\max}$ (gates/elem) & Ratio vs.\ Si \\
\midrule
Diamond   & $7.04 \times 10^{8}$ & $17.2\times$ \\
SiC (4H)  & $1.36 \times 10^{8}$ & $3.3\times$ \\
GaN       & $8.02 \times 10^{7}$ & $2.0\times$ \\
Silicon   & $4.10 \times 10^{7}$ & $1.0\times$ \\
GaAs      & $2.49 \times 10^{7}$ & $0.6\times$ \\
\bottomrule
\end{tabular}
\end{table}

\subsection{Cooling Sweep}

Sweeps $h$ from natural air (10~W/m$^2\cdot$K) to microchannel
cooling (50{,}000~W/m$^2\cdot$K), revealing the diminishing-returns
regime where the conduction floor dominates. This makes visible a
critical phenomenon: beyond a material-dependent threshold,
improving the cooling infrastructure yields negligible temperature
reduction because the die's own thermal conductivity is the bottleneck.

\subsection{Paradigm Comparison}

Compares CMOS, adiabatic, and reversible paradigms head-to-head at
the same configuration. At 7~nm and 1~GHz, adiabatic switching
consumes $1.22 \times 10^{-19}$~J/gate versus CMOS's $2.45 \times
10^{-16}$~J/gate --- a $2{,}006\times$ energy reduction. The crossover
frequency where CMOS becomes competitive is $f_c = 2.0$~THz,
well above practical operating frequencies.

\subsection{Thermal Headroom Map}

Analyses per-block thermal budget utilisation on a heterogeneous SoC.
For each functional block, it computes the steady-state temperature,
thermal headroom ($T_{\mathrm{limit}} - T_{\max}$), whether the block
is the thermal bottleneck, and how much additional density it could
accommodate. A representative modern SoC (CPU cluster, GPU array, SRAM
cache, I/O controller) typically reveals that cache and I/O blocks
utilise less than 5\% of their thermal budget while the CPU cluster is
at the limit.

\subsection{Power Redistribution Optimizer}

Given a total power budget $P_{\mathrm{budget}}$ and thermal limit
$T_{\mathrm{limit}}$, distributes gate density across functional
blocks to maximise total throughput:
\begin{equation}
  \max \sum_i D_i \cdot N_i \cdot a_i \cdot f \qquad
  \text{s.t.} \quad T_i \leq T_{\mathrm{limit}}, \;
  \sum_i P_i \leq P_{\mathrm{budget}}
  \label{eq:optim}
\end{equation}
The solver identifies whether the design is \emph{thermally limited}
(cannot use all power budget) or \emph{power limited} (thermal headroom
available). On a representative SoC, this yields up to $1.9\times$
throughput improvement over the unoptimised layout.

\subsection{Full Design Exploration}

A single-call entry point that runs all of the above analyses and
produces a comprehensive design space report with automatically
generated insights.


% ── 4. Technology Roadmap ────────────────────────────────────────────────
\section{Technology Roadmap Projection}
\label{sec:roadmap}

The \texttt{TechnologyRoadmap} module projects four metrics across 10
technology nodes (130~nm to 1.4~nm):

\begin{enumerate}[nosep]
  \item \textbf{Energy per gate switch} for each paradigm, showing how
    CMOS dynamic energy scales as $\sim V_{dd}^2 \cdot C_L$ while
    adiabatic energy scales as $\sim f$.
  \item \textbf{Landauer gap} (ratio of actual energy to $k_B T \ln 2$),
    showing convergence toward the thermodynamic limit at advanced
    nodes.
  \item \textbf{Paradigm crossover frequency} --- the frequency below
    which adiabatic logic consumes less energy than CMOS at each node.
  \item \textbf{Thermal wall} --- the maximum gate density sustainable
    on each substrate at each node, accounting for both energy model
    scaling and material thermal limits.
\end{enumerate}

Key finding: the Landauer gap for CMOS monotonically decreases from
$\sim 10^7\times$ at 130~nm to $\sim 10^4\times$ at 1.4~nm, but
remains orders of magnitude above the limit --- confirming that room
for improvement exists even at the most advanced nodes.


% ── 5. Validation ───────────────────────────────────────────────────────
\section{Validation}
\label{sec:validation}

Aethermor includes a self-contained validation suite (133 checks,
executable with a single command) that verifies every physics model
against published reference data. We organize validation into five
categories.

\subsection{Constants and Material Properties}

Fundamental constants are compared to CODATA~2018 exact
values~\cite{codata2018} with $< 0.01\%$ tolerance. Material properties
(thermal conductivity, density, specific heat, bandgap) are compared to
CRC Handbook~\cite{crc97} reference values with $< 0.1\%$ tolerance.
Derived quantities (thermal diffusivity $\alpha = k / \rho c_p$,
volumetric heat capacity) are verified for self-consistency.

\subsection{Energy Model Calibration}

CMOS supply voltage ($V_{dd}$) at five reference nodes is compared to
ITRS~2013~\cite{itrs2013} and IRDS~2022~\cite{irds2022} published
values, with all errors below 3\%. Key physics constraints are verified:
CMOS energy exceeds the Landauer limit (second law compliance), leakage
energy increases with temperature (Arrhenius), and adiabatic energy is
below CMOS energy at low frequency.

\subsection{Thermal Solver Verification}

The 3D Fourier solver is verified against the analytical solution for a
uniformly heated slab with convective cooling~\cite{carslaw1959}. The
3D solution is bounded between ambient temperature and the
(conservative) 1D analytical estimate. Energy conservation
($|E_{\mathrm{in}} - E_{\mathrm{out}} - \Delta U| / E_{\mathrm{in}}$)
is verified to be below 5\%, with observed error of 0.00\% in all runs.

\subsection{Inverse Design Consistency}

Round-trip consistency is verified for all inverse design tools:
\texttt{find\_max\_density} gives $T \approx T_{\mathrm{limit}}$ at
$D_{\max}$; \texttt{find\_min\_cooling} gives $T \approx
T_{\mathrm{limit}}$ at $h_{\min}$; the power redistribution optimizer
satisfies its budget and thermal constraints; the headroom map obeys
$T + \text{headroom} = T_{\mathrm{limit}}$ for every block.

\subsection{Reproducibility}

All functions that use the 3D binary-search solver are verified to
return identical results across independent runs with the same inputs.
The entire test suite (308~tests) and validation suite (133~checks)
execute deterministically with seeded random number generators.

\begin{table}[t]
\centering
\caption{Validation summary (133 checks total).}
\label{tab:validation}
\small
\begin{tabular}{@{}llrl@{}}
\toprule
Category & Reference & Checks & Status \\
\midrule
Fundamental constants     & CODATA 2018         & 6  & Pass \\
Landauer limit            & Landauer (1961)     & 5  & Pass \\
Material properties       & CRC Handbook        & 18 & Pass \\
CMOS energy model         & ITRS/IRDS           & 13 & Pass \\
Fourier solver            & Carslaw \& Jaeger   & 5  & Pass \\
Analytical 1D model       & Manual R-model      & 7  & Pass \\
Max density reciprocity   & 3D $\leftrightarrow$ analytical & 5 & Pass \\
Min cooling inverse       & Constraint round-trip & 4 & Pass \\
Optimizer constraints     & Budget/thermal/binding & 9 & Pass \\
Headroom map physics      & $T + h = T_{\mathrm{limit}}$ & 11 & Pass \\
Cooling stack resistance  & Incropera \& DeWitt & 4  & Pass \\
Tech roadmap monotonicity & 10 nodes, gap $> 1$ & 28 & Pass \\
Dimensional analysis      & Unit consistency    & 4  & Pass \\
Full exploration          & Response schema     & 11 & Pass \\
Reproducibility           & Deterministic       & 3  & Pass \\
\bottomrule
\end{tabular}
\end{table}


% ── 6. Representative Results ────────────────────────────────────────────
\section{Representative Results}
\label{sec:results}

We present three results that illustrate Aethermor's capabilities and
the physical insights they enable.

\subsection{Material Selection: Substrate Determines Density Ceiling}

At 7~nm, 1~GHz, with forced-air cooling ($h = 1{,}000$~W/m$^2\cdot$K),
diamond sustains $7.0 \times 10^8$~gates/element --- $17\times$
more than silicon ($4.1 \times 10^7$) and $28\times$ more than GaAs
($2.5 \times 10^7$). This quantifies a well-known qualitative
principle: substrate thermal conductivity sets the density ceiling. The
advantage persists across all cooling levels because the conduction
floor (determined by $k$) is the binding constraint at high $h$.

\subsection{Paradigm Crossover: Adiabatic Advantage at Practical Frequencies}

At 7~nm and room temperature, the CMOS $\leftrightarrow$ adiabatic
crossover occurs at $f_c = 2.0$~THz. Below this frequency --- which
includes the entire practical range for digital logic --- adiabatic
switching is more energy-efficient. The advantage at 1~GHz is
$2{,}006\times$. This holds across all technology nodes, with $f_c$
varying from $\sim$800~GHz at 130~nm to $\sim$3~THz at 1.4~nm.

\subsection{Heterogeneous SoC: Thermal Budget Redistribution}

On a representative modern SoC (CPU cluster at 5~nm, GPU array at
7~nm, SRAM cache at 7~nm, I/O controller at 14~nm), the headroom map
reveals that the I/O and cache blocks use less than 5\% of their
thermal budget. The power redistribution optimizer identifies that
reallocating compute density from underutilised blocks toward the
CPU cluster yields up to $1.9\times$ total throughput improvement ---
without changing the cooling architecture or exceeding any block's
thermal limit.


% ── 7. Comparison with Existing Tools ────────────────────────────────────
\section{Comparison with Existing Tools}
\label{sec:comparison}

Table~\ref{tab:comparison} compares Aethermor with existing thermal
analysis tools. The key differentiator is the inverse design
capability: Aethermor is the only open-source tool that can find
optimal configurations rather than merely evaluating given ones.

\begin{table}[t]
\centering
\caption{Feature comparison with existing tools.}
\label{tab:comparison}
\small
\begin{tabular}{@{}lccc@{}}
\toprule
Feature & Aethermor & HotSpot & COMSOL \\
\midrule
Forward thermal sim      & 3D Fourier & Compact & FEM \\
Inverse design           & 8 tools    & ---     & Script \\
Landauer gap tracking    & Yes        & ---     & --- \\
Multi-paradigm           & 4          & ---     & --- \\
Multi-material           & 9          & Si      & Yes \\
Cooling stack            & Multi-layer & Partial & Yes \\
Tech roadmap             & 10 nodes   & ---     & --- \\
SoC headroom map         & Per-block  & ---     & --- \\
Power redistribution     & Yes        & ---     & --- \\
Interactive UI           & Browser    & ---     & Desktop \\
Open source              & Apache 2.0 & BSD     & --- \\
Cost                     & Free       & Free    & \$25k+/yr \\
\bottomrule
\end{tabular}
\end{table}


% ── 8. Limitations and Future Work ──────────────────────────────────────
\section{Limitations and Future Work}
\label{sec:limitations}

We document limitations with the same rigour as capabilities.

\paragraph{Isotropic thermal model.}
All materials are treated as isotropic. Graphene's extreme in-plane
versus cross-plane anisotropy ($5{,}000$ vs.\ $\sim$6~W/m$\cdot$K)
is approximated with a single effective value. Extending
\texttt{FourierThermalTransport} to tensor conductivity is a priority
for future work.

\paragraph{No interconnect power.}
Energy models account for gate switching but not wire dissipation,
which is a significant fraction of total power at advanced nodes.

\paragraph{Simplified leakage.}
CMOS leakage uses a constant reference current ($I_0 = 1$~nA/gate)
across all nodes, ignoring short-channel effects (DIBL, gate tunnelling)
that increase leakage at smaller nodes. At room temperature, leakage is
a small fraction of dynamic energy, so this has minimal impact on
paradigm comparisons.

\paragraph{Validation scope.}
The thermal model has been validated against published specifications for four
real chips (NVIDIA A100, Apple M1, AMD EPYC 9654, Intel i9-13900K) and
produces correct-order-of-magnitude junction temperature predictions
(33 checks completed). Die-level correlation with proprietary floorplan
data or direct infrared thermal imaging is a planned next step.

\paragraph{Steady-state focus.}
The Fourier solver supports transient evolution but the inverse design
tools target steady-state. Modelling dynamic workloads, power gating
transients, and DVFS is future work.

\paragraph{1D cooling stack.}
The \texttt{CoolingStack} does not model 2D/3D spreading resistance
or temperature-dependent thermal conductivity. For package-level
analysis, these effects can be significant and will be addressed in a
future release.


% ── 9. Availability and Reproducibility ─────────────────────────────────
\section{Availability and Reproducibility}
\label{sec:availability}

Aethermor is released under the Apache~2.0 license at:

\begin{center}
\url{https://github.com/Yoder23/aethermor}
\end{center}

\noindent To reproduce all results in this paper:

\begin{lstlisting}
git clone https://github.com/Yoder23/aethermor
cd aethermor
pip install -e ".[all]"
python -m aethermor.validation.validate_all   # 133 checks
python -m pytest tests/ -v          # 308 tests
aethermor dashboard                 # Interactive UI
\end{lstlisting}

\noindent All random number generators are seeded for deterministic
output. The validation suite encodes the exact numerical checks
presented in this paper and can be run on any machine with
Python~$\geq$~3.10 and NumPy.


% ── 10. Conclusion ──────────────────────────────────────────────────────
\section{Conclusion}
\label{sec:conclusion}

Aethermor provides an integrated open-source toolkit for inverse thermal
design in thermodynamic computing research.  By combining calibrated
energy models, validated thermal solvers, and systematic inverse
design algorithms, it compresses the exploratory workflow for material
selections, cooling trade-offs, paradigm crossovers, and
heterogeneous SoC optimisations into interactive, sub-second queries. The
built-in validation suite of 133 cross-checks against published data
ensures that every result is grounded in verified physics. As
transistor scaling pushes designs closer to the Landauer limit, tools
like Aethermor become essential for navigating the narrowing design
space between what physics allows and what engineering can achieve.


% ── References ──────────────────────────────────────────────────────────
\begin{thebibliography}{20}

\bibitem{landauer1961}
R.~Landauer,
``Irreversibility and heat generation in the computing process,''
\emph{IBM J.\ Res.\ Dev.}, vol.~5, no.~3, pp.~183--191, 1961.

\bibitem{dennard1974}
R.~H. Dennard, F.~H. Gaensslen, H.~N. Yu, V.~L. Rideout, E.~Bassous,
and A.~R. LeBlanc,
``Design of ion-implanted MOSFET's with very small physical dimensions,''
\emph{IEEE J.\ Solid-State Circuits}, vol.~9, no.~5, pp.~256--268, 1974.

\bibitem{esmaeilzadeh2011}
H.~Esmaeilzadeh, E.~Blem, R.~St.~Amant, K.~Sankaralingam, and
D.~Burger,
``Dark silicon and the end of multicore scaling,''
in \emph{Proc.\ ISCA}, 2011, pp.~365--376.

\bibitem{skadron2004}
K.~Skadron, M.~R. Stan, K.~Sankaranarayanan, W.~Huang, S.~Velusamy,
and D.~Tarjan,
``Temperature-aware microarchitecture: Modeling and implementation,''
\emph{ACM Trans.\ Archit.\ Code Optim.}, vol.~1, no.~1, pp.~94--125,
2004.

\bibitem{itrs2013}
International Technology Roadmap for Semiconductors,
\emph{ITRS 2013 Edition}, Semiconductor Industry Association, 2013.

\bibitem{irds2022}
IEEE International Roadmap for Devices and Systems,
\emph{IRDS 2022 Edition}, IEEE, 2022.

\bibitem{codata2018}
E.~Tiesinga, P.~J. Mohr, D.~B. Newell, and B.~N. Taylor,
``CODATA recommended values of the fundamental physical constants: 2018,''
\emph{Rev.\ Mod.\ Phys.}, vol.~93, no.~2, 025010, 2021.

\bibitem{crc97}
W.~M. Haynes, Ed.,
\emph{CRC Handbook of Chemistry and Physics}, 97th ed.,
CRC Press, 2016.

\bibitem{carslaw1959}
H.~S. Carslaw and J.~C. Jaeger,
\emph{Conduction of Heat in Solids}, 2nd ed.,
Oxford University Press, 1959.

\bibitem{morkoc2006}
H.~Morko\c{c},
``SiC and GaN based devices --- review and challenges,''
in \emph{Proc.\ SPIE}, vol.~6127, 2006.

\bibitem{morkoc2008}
H.~Morko\c{c},
\emph{Handbook of Nitride Semiconductors and Devices},
Wiley, 2008.

\bibitem{incropera2011}
F.~P. Incropera, D.~P. DeWitt, T.~L. Bergman, and A.~S. Lavine,
\emph{Fundamentals of Heat and Mass Transfer}, 7th ed.,
Wiley, 2011.

\end{thebibliography}

\end{document}
